\chapter[Lezione 2 --- 30 settembre 2026]{Lezione 2: numeri interi, razionali e fattoriale}
\label{chap:lezione-02}

\section{Numeri interi}
\label{sec:lez02-numeri-interi}

Dato un insieme numerico \(A\), si indica con \(-A\) l'insieme degli
opposti dei suoi elementi:
\[
  -A:=\{-a:a\in A\}.
\]
Poiché \(0\in\N\), l'insieme \(-\N\) contiene lo zero e gli opposti
dei naturali positivi. L'unione di \(\N\) e \(-\N\) permette quindi
di introdurre i numeri interi.

\begin{definition}[Numeri interi]
\label{def:numeri-interi}
L'insieme dei \emph{numeri interi} è
\[
  \Z:=\N\cup(-\N).
\]
\end{definition}

Si assume l'\emph{assioma di esistenza e unicità dell'opposto}: per ogni
\(m\in\Z\) esiste un unico \(n\in\Z\) tale che
\[
  m+n=0.
\]
Tale elemento si indica con \(-m\) ed è detto \emph{opposto} di \(m\).

L'ordinamento dei naturali, introdotto nella
\cref{sec:lez01-ordine-naturali}, si estende agli interi. La
moltiplicazione per un intero positivo conserva il verso delle
disuguaglianze, mentre la moltiplicazione per un intero negativo
lo inverte.

\begin{proposition}[Compatibilità dell'ordine in \(\Z\)]
\label{prop:compatibilita-ordine-interi}
Per ogni \(a,b,c\in\Z\), se \(a\le b\), allora
\[
  c>0\Longrightarrow ac\le bc,
  \qquad
  c<0\Longrightarrow ac\ge bc.
\]
\end{proposition}

Per \(c=0\) si ha invece \(ac=bc=0\). La dimostrazione della
\cref{prop:compatibilita-ordine-interi} non viene sviluppata
in questa sede.

Si consideri ora il problema di determinare, per \(a\in\Z\), un
intero \(b\) tale che
\[
  ab=1.
\]
Quando esiste, tale elemento è detto \emph{inverso moltiplicativo}
di \(a\) e si indica con \(a^{-1}\). Negli interi, tuttavia, non
ogni elemento ammette un inverso moltiplicativo.

\begin{proposition}[Inversi moltiplicativi in \(\Z\)]
\label{prop:inversi-interi}
Gli unici interi che ammettono un inverso moltiplicativo in
\(\Z\) sono \(-1\) e \(1\). In particolare,
\[
  (-1)^{-1}=-1,
  \qquad
  1^{-1}=1.
\]
\end{proposition}

\begin{proof}
Siano \(a,b\in\Z\) tali che \(ab=1\). Il prodotto è positivo,
quindi \(a\) e \(b\) hanno lo stesso segno.

Se entrambi sono positivi, si ha \(a,b\ge1\). Poiché \(ab=1\),
necessariamente \(a=b=1\). Se entrambi sono negativi, gli interi
positivi \(-a\) e \(-b\) soddisfano \((-a)(-b)=1\), da cui
\(a=b=-1\).

Viceversa, \(1\cdot1=1\) e \((-1)(-1)=1\), quindi entrambi gli
interi indicati ammettono un inverso moltiplicativo in \(\Z\).
\end{proof}

\section{Numeri razionali}
\label{sec:lez02-numeri-razionali}

La \cref{prop:inversi-interi} mostra che in \(\Z\) soltanto
\(-1\) e \(1\) ammettono un inverso moltiplicativo. Si introduce
pertanto un insieme numerico più ampio, nel quale sia possibile
rappresentare i rapporti tra interi con denominatore non nullo.

Per due insiemi \(A\) e \(B\), la \emph{differenza insiemistica}
è definita da
\[
  A\setminus B:=\{x\in A:x\notin B\}.
\]
In particolare, \(\N\setminus\{0\}\) è l'insieme dei naturali
strettamente positivi.

\begin{definition}[Numeri razionali]
\label{def:numeri-razionali}
L'insieme dei \emph{numeri razionali} è
\[
  \Q:=\left\{\frac{m}{n}:m\in\Z,\ n\in\N\setminus\{0\}\right\}.
\]
\end{definition}

Si ottiene così la catena di inclusioni
\[
  \N\subseteq\Z\subseteq\Q.
\]

Il denominatore è scelto positivo, riportando l'eventuale segno
negativo al numeratore. Uno stesso razionale può essere rappresentato
da frazioni diverse, per esempio \(\frac12=\frac24\). Riducendo
una frazione ai minimi termini si ottiene una rappresentazione
\[
  \frac{m}{n},\qquad \operatorname{mcd}(m,n)=1,
\]
dove il massimo comune divisore è quello positivo di \(|m|\) e \(n\).
Con \(n>0\), la rappresentazione ridotta è unica; in particolare,
lo zero si rappresenta mediante \(\frac01\).

In \(\Q\) si assume l'\emph{assioma dell'inverso moltiplicativo}:
\[
  \forall a\in\Q\setminus\{0\}\;\exists!\,b\in\Q:
  \quad ab=1.
\]
L'unico elemento \(b\) che soddisfa questa condizione si indica
con \(a^{-1}\) ed è detto anche \emph{reciproco} di \(a\). Per esempio,
\[
  2^{-1}=\frac12.
\]

\section[Irrazionalità della radice quadrata di 2]{Irrazionalità di \(\sqrt{2}\)}
\label{sec:lez02-irrazionalita-radice-due}

L'insieme dei razionali non è sufficiente a risolvere tutte le
equazioni polinomiali. Si consideri, per esempio, il polinomio
\[
  p(x)=x^2-2,
  \qquad p\in\Q[x],
\]
dove \(\Q[x]\) indica l'insieme dei polinomi nella variabile
\(x\) a coefficienti razionali. Si dimostrerà che l'equazione
\(x^2-2=0\) non ammette soluzioni razionali. Il risultato sarà
anche espresso nella forma \(\sqrt{2}\notin\Q\): l'esistenza
della radice quadrata reale non negativa di \(2\) sarà giustificata
nella \cref{sec:lez03-completezza-radice-due}.

Si richiamano anzitutto le nozioni di numero pari e dispari,
necessarie per la dimostrazione.

\clearpage

\begin{definition}[Numeri pari e dispari]
\label{def:lez02-pari-dispari}
Un intero \(m\) si dice \emph{pari} se esiste \(k\in\Z\) tale che
\[
  m=2k,
\]
mentre si dice \emph{dispari} se esiste \(k\in\Z\) tale che
\[
  m=2k+1.
\]
\end{definition}

Ogni intero è pari oppure dispari, e le due possibilità sono
mutuamente esclusive.

\begin{lemma}[Parità del quadrato]
\label{lem:quadrato-pari}
Per ogni \(m\in\Z\), se \(m^2\) è pari, allora \(m\) è pari.
\end{lemma}

\begin{proof}
Si dimostra la contronominale, secondo lo schema introdotto nella
\cref{sec:lez01-insiemi-dimostrazioni}. Sia \(m\) dispari. Per la
\cref{def:lez02-pari-dispari}, esiste \(k\in\Z\) tale che
\[
  m=2k+1.
\]
Pertanto
\begin{align*}
  m^2
    &=(2k+1)^2\\
    &=4k^2+4k+1\\
    &=2\bigl(2k(k+1)\bigr)+1.
\end{align*}
Poiché \(2k(k+1)\in\Z\), il quadrato \(m^2\) è dispari.
La contronominale è verificata e il lemma segue.
\end{proof}

\begin{theorem}[Irrazionalità di \(\sqrt{2}\)]
\label{thm:irrazionalita-radice-due}
L'equazione \(x^2-2=0\) non ammette soluzioni in \(\Q\).
Equivalentemente, con la notazione introdotta sopra,
\[
  \sqrt{2}\notin\Q.
\]
\end{theorem}

\begin{proof}
Si procede per assurdo. Si supponga che esista \(x\in\Q\) tale che
\[
  x^2=2.
\]
Scrivendo \(x\) in forma ridotta, esistono \(m\in\Z\) e
\(n\in\N\setminus\{0\}\) tali che
\[
  x=\frac{m}{n},
  \qquad
  \operatorname{mcd}(m,n)=1.
\]
Sostituendo la rappresentazione di \(x\) nell'uguaglianza
\(x^2=2\), si ottiene
\[
  \left(\frac{m}{n}\right)^2=2.
\]
Moltiplicando per \(n^2>0\), segue
\[
  m^2=2n^2.
\]
Di conseguenza, \(m^2\) è pari e, per il \cref{lem:quadrato-pari},
anche \(m\) è pari. Esiste quindi \(h\in\Z\) tale che \(m=2h\).
Sostituendo nell'uguaglianza precedente,
\[
  4h^2=2n^2
  \quad\Longrightarrow\quad
  n^2=2h^2.
\]
Anche \(n^2\) è pari e, applicando nuovamente il
\cref{lem:quadrato-pari}, si conclude che \(n\) è pari.

Gli interi \(m\) e \(n\) sono dunque entrambi divisibili per \(2\),
in contraddizione con \(\operatorname{mcd}(m,n)=1\).
L'ipotesi iniziale è falsa: l'equazione \(x^2-2=0\) non ammette
soluzioni razionali.
\end{proof}

Si richiama ora la nozione di numero primo per enunciare
una generalizzazione del risultato precedente.

\begin{definition}[Numero primo]
\label{def:numero-primo}
Un numero naturale \(p>1\) si dice \emph{primo} se i suoi unici
divisori positivi sono \(1\) e \(p\).
\end{definition}

Per esempio, \(2\) e \(3\) sono primi, mentre \(4\) non lo è,
poiché è divisibile anche per \(2\). Il numero \(1\) non è primo,
come stabilito dalla definizione.

\begin{proposition}[Irrazionalità della radice di un numero primo]
\label{prop:irrazionalita-radice-primo}
Sia \(p\in\N\) un numero primo. L'equazione
\[
  x^2-p=0
\]
non ammette soluzioni in \(\Q\). Indicando con \(\sqrt{p}\) la
radice quadrata reale non negativa di \(p\), il risultato si esprime
nella forma
\[
  \sqrt{p}\notin\Q.
\]
\end{proposition}

La dimostrazione della \cref{prop:irrazionalita-radice-primo} non viene
sviluppata in questa sede.

\section{Fattoriale}
\label{sec:lez02-fattoriale}

\begin{definition}[Fattoriale]
\label{def:fattoriale}
Il \emph{fattoriale} di un numero naturale è definito ricorsivamente
dalle condizioni
\[
  0!:=1,
  \qquad
  (n+1)!:=(n+1)n!,
  \quad n\in\N.
\]
\end{definition}

Applicando la \cref{def:fattoriale}, si ottiene
\begin{align*}
  0!&=1,\\
  1!&=1\cdot0!=1,\\
  2!&=2\cdot1!=2,\\
  3!&=3\cdot2!=3\cdot2\cdot1.
\end{align*}
Per \(n\ge1\), il fattoriale ammette anche la rappresentazione
\[
  n!=1\cdot2\cdot\ldots\cdot n=\prod_{i=1}^{n}i.
\]

\begin{proposition}[Una disuguaglianza per il fattoriale]
\label{prop:fattoriale-maggiore-n}
Per ogni \(n\in\N\) vale
\[
  n!\ge n.
\]
\end{proposition}

\begin{proof}
Si consideri l'insieme
\[
  A:=\{n\in\N:n!\ge n\}.
\]
Per \(n=0\), dalla \cref{def:fattoriale} si ha
\[
  0!=1\ge0,
\]
quindi \(0\in A\).

Si procede per induzione su \(n\ge1\), applicando il
\cref{thm:principio-induzione} con valore iniziale \(p=1\).

\emph{Passo base.}
Per \(n=1\), si ha
\[
  1!=1\ge1.
\]

\emph{Passo induttivo.}
Sia \(n\ge1\) e si supponga che valga l'ipotesi induttiva
\[
  n!\ge n.
\]
Per la \cref{def:fattoriale} e la compatibilità dell'ordine
rispetto alla moltiplicazione per un numero positivo
(\cref{prop:compatibilita-ordine}), si ottiene
\[
  (n+1)!=(n+1)n!\ge(n+1)n\ge n+1,
\]
poiché \(n+1>0\) e \(n\ge1\). Il passo induttivo è quindi
verificato.

Per il principio di induzione, la disuguaglianza vale per ogni
\(n\ge1\). Poiché è verificata anche per \(n=0\), si conclude
che \(A=\N\), e dunque la tesi.
\end{proof}
